\clearpage
\section{Chapter 2: Hubble Law \& Cosmological Redshift}

\subsection{Questions}

\begin{enumerate}[label=2.\arabic*]
  \item \label{q:2.1} \textbf{Redshift and Scale Factor}

  A galaxy is observed at redshift $z=2.5$.
  \begin{enumerate}[label=(\alph*)]
    \item Find the scale factor at emission, $a(t_{\mathrm{emit}})$, assuming $a_0=1$ today.
    \item If the emitted wavelength is $\lambda_{\mathrm{emit}}=121.6\,\mathrm{nm}$ (Lyman-$\alpha$), find $\lambda_{\mathrm{obs}}$.
    \item If its comoving distance is $\chi=1.2\,\mathrm{Gpc}$, compute the proper distance at emission and at present.
  \end{enumerate}

  \item \label{q:2.2} \textbf{Hubble Law in Practice}

  Take $H_0 = 67.66\,\mathrm{km\,s^{-1}\,Mpc^{-1}}$.
  \begin{enumerate}[label=(\alph*)]
    \item Convert $H_0$ to SI units ($\mathrm{s^{-1}}$), using $1\,\mathrm{Mpc}=3.086\times 10^{22}\,\mathrm{m}$.
    \item Estimate the Hubble time $t_H = 1/H_0$ in Gyr.
    \item For a nearby galaxy at $D=100\,\mathrm{Mpc}$, estimate recession speed $v$ from Hubble's law, and then estimate $z\approx v/c$ (low-$z$ limit).
  \end{enumerate}

  \item \label{q:2.3} \textbf{Hubble Tension as a Significance Test}

  Two measurements are reported:
  \[
    H_0^{\mathrm{SH0ES}} = 73.04 \pm 1.04,\qquad
    H_0^{\mathrm{Planck}} = 67.66 \pm 0.42
    \quad (\mathrm{km\,s^{-1}\,Mpc^{-1}}).
  \]
  \begin{enumerate}[label=(\alph*)]
    \item Compute the absolute difference $\Delta H_0$.
    \item \textit{(Hint: use independent Gaussian errors)} Compute combined uncertainty
    $\sigma_{\Delta} = \sqrt{\sigma_1^2+\sigma_2^2}$.
    \item Compute significance $N_\sigma = \Delta H_0/\sigma_\Delta$ and interpret briefly.
  \end{enumerate}

  \item \label{q:2.4} \textbf{Luminosity Distance and Flux}

  In a flat universe ($k=0$), use
  \[
    d_L=(1+z)\chi,
    \qquad
    F=\frac{L}{4\pi d_L^2}.
  \]
  A source is observed at $z=2.5$ with comoving distance $\chi=1.2\,\mathrm{Gpc}$ and intrinsic luminosity
  \[
    L=1.0\times 10^{37}\,\mathrm{W}.
  \]
  \begin{enumerate}[label=(\alph*)]
    \item Compute $d_L$ in Gpc.
    \item Convert $d_L$ to meters and compute the observed flux $F$.
    \item Using $a_0=1$, show that $d_L=(1+z)^2D_{\mathrm{proper,emit}}$ and verify numerically with $D_{\mathrm{proper,emit}}=a_{\mathrm{emit}}\chi$.
  \end{enumerate}
\end{enumerate}

\bigskip
\bigskip
\bigskip

\subsection{Solutions}

\subsubsection*{Solution 2.1: Redshift and Scale Factor}

\begin{enumerate}[label=(\alph*)]
  \item Using
  \[
    z = \frac{1}{a_{\mathrm{emit}}}-1,
  \]
  we get
  \[
    a_{\mathrm{emit}} = \frac{1}{1+z} = \frac{1}{3.5} \approx 0.286.
  \]

  \item Cosmological redshift gives
  \[
    \lambda_{\mathrm{obs}} = (1+z)\lambda_{\mathrm{emit}} = 3.5\times 121.6\,\mathrm{nm}
    = 425.6\,\mathrm{nm}.
  \]

  \item Proper distance is $D_{\mathrm{proper}}(t)=a(t)\chi$:
  \[
    D_{\mathrm{proper,emit}} = a_{\mathrm{emit}}\chi \approx 0.286\times 1.2 = 0.343\,\mathrm{Gpc},
  \]
  \[
    D_{\mathrm{proper,0}} = a_0 \chi = 1\times 1.2 = 1.2\,\mathrm{Gpc}.
  \]
\end{enumerate}

\subsubsection*{Solution 2.2: Hubble Law in Practice}

\begin{enumerate}[label=(\alph*)]
  \item
  \[
    H_0 = 67.66\,\frac{\mathrm{km}}{\mathrm{s\,Mpc}}
    = \frac{67.66\times 10^3\,\mathrm{m/s}}{3.086\times 10^{22}\,\mathrm{m}}
    \approx 2.19\times 10^{-18}\,\mathrm{s^{-1}}.
  \]

  \item
  \[
    t_H = \frac{1}{H_0} \approx \frac{1}{2.19\times 10^{-18}} = 4.56\times 10^{17}\,\mathrm{s}.
  \]
  With $1\,\mathrm{Gyr}=3.156\times 10^{16}\,\mathrm{s}$,
  \[
    t_H \approx 14.4\,\mathrm{Gyr}.
  \]

  \item
  \[
    v=H_0D = 67.66\times 100\,\mathrm{km/s} = 6766\,\mathrm{km/s}.
  \]
  For low redshift ($v\ll c$),
  \[
    z\approx \frac{v}{c} \approx \frac{6766}{3.0\times 10^5} \approx 2.26\times 10^{-2}.
  \]
  Since $v/c\approx 0.0226\ll 1$, the low-$z$ approximation is self-consistent.
\end{enumerate}

\subsubsection*{Solution 2.3: Hubble Tension as a Significance Test}

\begin{enumerate}[label=(\alph*)]
  \item
  \[
    \Delta H_0 = 73.04-67.66 = 5.38\,\mathrm{km\,s^{-1}\,Mpc^{-1}}.
  \]

  \item
  \[
    \sigma_\Delta = \sqrt{1.04^2+0.42^2} = \sqrt{1.258} \approx 1.12.
  \]

  \item
  \[
    N_\sigma = \frac{5.38}{1.12} \approx 4.8\sigma.
  \]
  A $\sim 4.8\sigma$ discrepancy is statistically significant, consistent with the notes' statement that this is the ``Hubble tension'' problem.
\end{enumerate}

\subsubsection*{Solution 2.4: Luminosity Distance and Flux}

\begin{enumerate}[label=(\alph*)]
  \item
  \[
    d_L=(1+z)\chi = 3.5\times 1.2\,\mathrm{Gpc}=4.2\,\mathrm{Gpc}.
  \]

  \item Using $1\,\mathrm{Gpc}=3.086\times10^{25}\,\mathrm{m}$:
  \[
    d_L=4.2\times 3.086\times10^{25}\,\mathrm{m}
    \approx 1.30\times10^{26}\,\mathrm{m}.
  \]
  Then
  \[
    F=\frac{L}{4\pi d_L^2}
    =\frac{1.0\times10^{37}}{4\pi(1.30\times10^{26})^2}
    \approx 4.7\times10^{-17}\,\mathrm{W\,m^{-2}}.
  \]

  \item Since $D_{\mathrm{proper,emit}}=a_{\mathrm{emit}}\chi$ and
  \[
    a_{\mathrm{emit}}=\frac{1}{1+z},
  \]
  we get
  \[
    D_{\mathrm{proper,emit}}=\frac{\chi}{1+z}
    \Rightarrow
    \chi=(1+z)D_{\mathrm{proper,emit}}.
  \]
  Therefore
  \[
    d_L=(1+z)\chi=(1+z)^2D_{\mathrm{proper,emit}}.
  \]
  Numerically, from Solution 2.1,
  \[
    D_{\mathrm{proper,emit}}\approx 0.343\,\mathrm{Gpc},
  \]
  so
  \[
    (1+z)^2D_{\mathrm{proper,emit}}=3.5^2\times0.343\approx 4.2\,\mathrm{Gpc},
  \]
  consistent with part (a).
\end{enumerate}


% ─────────────────────────────────────────────────────────────────────────────
